\section{What observations determine about memory, records, and time}
\label{op:collections}

The operational studies concern finite conventional models and ask a precise
question: which features of a model are fixed by a stated observation and
intervention menu?  The answers include positive identification theorems,
finite statistical designs, and explicit indistinguishable alternatives.
Their premises include ordinary probability or quantum theory, calibrated
settings, and a specified notion of state.  They do not constitute a
derivation of those primitives from a prephysical candidate.
Complete statements, proofs, and computational certificates appear in the
Mathematical Supplement under \emph{Observation, Memory, Control, and Time
in Finite Models}.

\subsection{Reduced dynamics require a predictive state}

For a finite controlled Markov model, an observation quotient has exact
autonomous dynamics for all starting distributions and all common
visible-history policies precisely when each action's pushed transition
row is constant within each observation fiber.  This is a controlled form
of strong lumpability \cite{op:wolfer}.  Uniform conditional row error
\(\delta\) gives a sharp complete-history total variation bound
\[
 1-(1-\delta)^N
\]
after \(N\) transitions.  An average stationary fit does not replace the
uniform premise.

A four-state stationary process illustrates the difference.  Its observed
singles are fair and every adjacent pair is independent and fair, yet its
triple law differs from the fair iid law by total variation \(1/4\).
Conversely, a hidden bit may alternate forever without affecting the
visible Markov law when its pushed rows coincide.  Hidden state is neither
automatically harmful nor automatically recoverable.
An operation on fine states can also distinguish two inputs that the
observation map identifies.  Such an operation is not determined by the
coarse state unless the intervention itself descends through the quotient.

\subsection{Memory depth, readout, and state count}

The parity-shift family has a uniform \(L\)-bit initial block and append
probabilities
\[
 P(b\mid x_1,\ldots,x_L)
 =\tfrac12[1+\eta(-1)^{x_1+\cdots+x_L+b}].
\]
It belongs to the established two-faced-chain setting
\cite{op:ryabko}.  Every contiguous block of length at most \(L\) has the
fair iid law, while the first \(L+1\) block differs by
\(|\eta|/2\) in total variation.  Therefore no common finite block depth
validates independence over unrestricted \(L\).

The parity innovations are iid under ideal readout because the entire
record probability factors in those coordinates.  Independent detector
bit flips alter that conclusion: consecutive observed parity windows
overlap in their detector bits, producing nonzero covariance through lag
\(L\).  Their count is generally insufficient for the ordered record.
For \(L=2,\eta=1/2,\epsilon=1/10\), a rule that attained more than \(95\%\)
power with \(40\) ideal innovations has only about \(49.26\%\) power after
the specified readout noise.  An exact finite recursion gives a minimum
of \(183\) innovations for the randomized, size-\(5\%\), count-only
experiment to attain \(95\%\) power.  Every shorter count horizon is
checked separately; this is not a minimum for all full-record procedures.

The joint matrix of adjacent \(L\)-bit past and future blocks has rank
\(2^L\) for nonzero \(\eta\).  The explicit shift-state model attains that
number, so it is the exact minimum for classical transition-emitting
realizations of this process.  Independent symmetric readout preserves
the rank when its flip probability is below \(1/2\), while potentially
making the small singular values difficult to estimate.
This model differs from a stationary random-phase concatenation of
perturbed \(k\)-bit blocks; the latter's separate \(2k-1\) realization
result does not describe the parity-shift family.

Population rank is not by itself a finite-data certificate.  For \(n\)
iid complete records distributed over \(K\) joint-matrix cells, the
empirical matrix obeys the confidence radius
\[
 \|\widehat H-H\|_F\le
 \frac{\sqrt{1-1/K}+\sqrt{\log(1/\alpha)}}{\sqrt n}
\]
with probability at least \(1-\alpha\).
Combining this radius with the singular tail produces rank tests and
confidence lower bounds on total variation from every normalized
low-rank competitor.  Exact matching of competitor marginals strengthens
the population bound; that restriction must be justified, not silently
imposed.
For the specified four-output noisy example, the conservative \(5\%\)
size and \(95\%\) power designs require \(2976,5116,18213\) iid records
to reject state counts at most \(1,2,3\), respectively.
These are sufficient budgets for one particular bound, not optimal
sample sizes.  The independence requirement is across complete records;
dependence within a record is allowed.

\subsection{Sequential validity and power are different questions}

Against the complete fair iid observed law, a predeclared mixture of
coherent record likelihood ratios is a nonnegative martingale.
Crossing \(20\) gives an anytime false-rejection probability at most
\(5\%\), by the standard martingale argument \cite{op:howard}.
With weights \(1/[L(L+1)]\) over memory lengths, fixed
\(\eta=1/2,\epsilon=1/10\), and true \(L=2\), a component argument using
\(650\) disjoint three-bit parity windows proves power by \(1950\)
outputs of at least \(0.9519745696463602\ldots\).
This is a sufficient horizon for that alternative.  Every fixed
nondegenerate candidate with positive mixture weight is detected
eventually almost surely, but no finite horizon works uniformly over
unrestricted memory length.

Conditional-bias null envelopes permit dependence and require a
different interpretation.  A valid normalized forecast factor can have
negative long-run log growth even under its own correct alternative
forecast.  Clipping a Bernoulli forecast to the null interval gives an
explicit valid factor that maximizes one-step expected log growth,
consistent with established e-variable theory \cite{op:larsson}.
Enlarging an envelope can admit the memory-bearing alternative itself;
no test valid on that enlarged null can simultaneously promise high
rejection probability for the admitted process.

\subsection{Final-system records do not locate environmental structure}

Conditionally identical product fragments with overlap \(\gamma\) give
visibility \(|\gamma|^N\).  For real \(0\le\gamma\le1\), the final binary
family \([1+\gamma^N\cos\phi]/2\) also has an elementary two-state classical
realization with a setting-dependent readout.  This reproduces that
specified binary family; it makes no claim to simulate arbitrary quantum
experiments.

Even complete final-system channel knowledge does not identify
tensor-product factorization of inaccessible conditional fragment
states.  A common environment-only controlled-phase unitary preserves
the full reduced channel, including retained ancillas, while changing a
product branch into a nonproduct branch across the fixed fragment
partition.  Complete later histories remain equal only when those
transformed fragments are never accessed or recoupled.
The distinction concerns quantum conditional states, not necessarily
the classical distribution of a particular fragment measurement.

With branch preparation and local fragment access added, the explicit
pair becomes distinguishable.  At overlap \(\gamma=3/5\), a calibrated
local \(X\) readout has binary probability gap \(384/625\).
The pair is still separated if each output law has total variation error
at most \(\varepsilon\) and \(384/625-2\varepsilon>0\).
Thus a change in access assumptions changes the operational conclusion.

\subsection{Exact controls and uncertain controls}

Strictly positive, linearly minimal two-state Mealy models can have
identical passive output laws without being state-label permutations
\cite{op:vanluyten}.  In the explicit pair, a common silent control
\[
 K(a,b)=\begin{pmatrix}1-a&a\\b&1-b\end{pmatrix}
\]
separates the next-zero probability by \((2a-b)/35\).
All controlled laws agree precisely when \(b=2a\).
A reset can separate this pair yet leave other minimal partners, and
one fixed silent control leaves a local continuum throughout the
unrestricted strictly positive minimal class with unanchored preparation.

There is a sharp positive two-state result.  Consider stationary,
binary-output, strictly positive transition-emitting models with
linear minimal dimension two and a preparation that is not separately
anchored.  Two exactly known controls \(K_j=K(a_j,b_j)\), in one common
labelled hidden basis and fixed across the compared models, identify
the numerical realization from controlled population laws if and only if
\[
 a_1b_2-a_2b_1\ne0.
\]
At most one control followed by two outputs per setting suffices for an
explicit reconstruction, using the same starting preparation across
settings.  The result requires exact common controls and exact
probabilities; it does not give a finite-sample guarantee or a uniform
stability bound.

In arbitrary finite dimension, an exactly labelled cycle and reset
prepare every pure state; observability then gives reconstruction from
records with at most \(n\) outputs.

Calibration uncertainty reopens equivalence.  For an explicit positive
pair of control matrices, fixing one exactly while allowing the other
a full row-total-variation ball of radius \(\varepsilon\) leaves a
contained first-preparation-coordinate interval
\[
 \left[\frac49,\frac49+
       \min\left\{\frac1{15},\frac{8\varepsilon}{9}\right\}\right]
\]
with identical complete controlled laws.  It is a lower bound on an
identified set, not its full characterization.  In a separate fixed-pair
reset experiment, uniform separation fails at a \(20\%\) preparation
nuisance threshold.  Perfect iid reset trials for that pair have an
exact randomized fixed-horizon minimum of \(468\) trials at size \(5\%\)
and power at least \(95\%\).  These numbers concern different
experiments and are not tolerances or sample budgets for the
two-control reconstruction.

\subsection{Sampling time can hide irreversible circulation}

Two explicitly positive three-state continuous-time generators have
the same full transition matrix at every nonnegative integer time.
One is reversible and one has nonzero stationary cycle current.
Their sampled equivalence survives every common adaptive policy whose
controls and readouts, including silent state maps, all occur on the
integer grid while the generator remains fixed between boundaries.
This does not include off-grid controls, persistent drives, or
continuous monitoring.
An anchored state-one preparation and a calibrated quarter-time readout
resolving states two and three gives the explicit distinguishing
contrast \(2e^{-3}/\sqrt3\).

Full exact transition matrices at two positive times of irrational ratio
uniquely determine a finite generator in a common basis.
Nevertheless, rates tending to infinity can make finitely many sampled
kernels arbitrarily close while the generators remain a fixed distance
apart.  Exact uniqueness and uniform inverse stability therefore differ.
The compact, bounded-parameter case admits qualitative uniform
continuity, but no general quantitative finite-data clock-calibration
theorem follows here.  These are results within the classical
matrix-exponential inverse problem \cite{op:baake,op:yue}; they
presuppose a physical time assignment.
